\section{Resultados Potenciais e Hipóteses}

\begin{frame}{Resultados Potenciais (Potential Outcomes)}
    \begin{block}{O que aconteceria se...}
        Antes da divisão do sorteio, cada paciente possui dois resultados potenciais:
        \begin{enumerate}
            \item O resultado obtido caso fosse sorteado para o grupo de tratamento.
            \item O resultado obtido caso fosse sorteado para o grupo de controle.
        \end{enumerate}
    \end{block}
    \begin{block}{A Metáfora do Bilhete de Dois Lados}
        Imagine que cada um dos $31$ pacientes possui um bilhete. O lado esquerdo é a sua melhora com BTA, e o lado direito é sua melhora com salina.
        \begin{itemize}
            \item Para os 15 sorteados no tratamento, revelamos a face esquerda (BTA).
            \item Para os 16 sorteados no controle, revelamos a face direita (salina).
            \item A outra metade de cada bilhete permanece eternamente \textbf{desconhecida}.
        \end{itemize}
    \end{block}
    \centering
    \textit{Nosso desafio estatístico é comparar duas distribuições tendo acesso apenas a metade dos dados de cada bilhete.}
\end{frame}

\begin{frame}{As Hipóteses do Teste Clínico}
    \begin{block}{Hipótese Nula ($H_0$): O Modelo do Acaso}
        A distribuição de todos os $31$ resultados potenciais sob o tratamento é idêntica à distribuição sob o controle. A Toxina Botulínica Tipo A (BTA) não difere do soro fisiológico.
        \begin{itemize}
            \item Sob $H_0$, as faces esquerda e direita do bilhete de cada paciente têm o mesmo valor.
            \item Qualquer diferença observada deve-se puramente ao acaso do sorteio dos grupos.
        \end{itemize}
    \end{block}
    \begin{block}{Hipótese Alternativa ($H_1$): A Presença de Efeito}
        A distribuição de resultados potenciais sob o tratamento é diferente daquela sob o controle. O remédio realmente causa um efeito diferente de um placebo.
    \end{block}
\end{frame}
